\documentclass[10pt,a4paper]{article}

\usepackage[T1]{fontenc}
\usepackage[utf8]{inputenc}
\usepackage[a4paper,margin=19mm]{geometry}
\usepackage{amsmath}
\usepackage{array}
\usepackage{enumitem}
\usepackage{listings}
\usepackage{url}
\usepackage{microtype}

\setlist{nosep,leftmargin=*}
\setlength{\parindent}{1em}
\setlength{\parskip}{0pt}
\setlength{\textfloatsep}{6pt plus 2pt minus 2pt}
\setlength{\intextsep}{6pt plus 2pt minus 2pt}
\setlength{\abovedisplayskip}{5pt plus 2pt minus 2pt}
\setlength{\belowdisplayskip}{5pt plus 2pt minus 2pt}
\emergencystretch=2em
\widowpenalty=10000
\clubpenalty=10000
\Urlmuskip=0mu plus 1mu

\lstset{
  basicstyle=\ttfamily\small,
  columns=fullflexible,
  keepspaces=true,
  showstringspaces=false,
  frame=single,
  aboveskip=5pt,
  belowskip=5pt,
  breaklines=true
}

\newcommand{\meet}{\mathbin{\sqcap}}
\newcommand{\seq}{\mathbin{\cdot}}

\title{Static Type Checking for Forth by Symbolic Stack Effects}
\author{Jaanus Pöial\\Tallinn University of Technology}
\date{}

\begin{document}
\maketitle
\vspace{-1.5em}

\begin{abstract}
This paper describes a source-level static checker intended for Forth systems
whose vocabulary and parsing rules can be described externally.  The checker
reads a type hierarchy, stack-effect and parser specifications, and Forth
source.  It represents each phrase by a pair of symbolic stacks.  Indexed
symbols retain the identity of cells moved or copied by words such as
\texttt{DUP}, \texttt{OVER}, and \texttt{SWAP}; subtype-aware unification
checks adjacent effects and propagates the most precise available type.
Alternative paths are reduced to a common stack contract.  Repetition uses a
finite one-execution versus two-execution test rather than a general loop
invariant calculation.  Parser words, defining words, literals, and control
structures are profile data, so checking is not tied to one fixed Forth
syntax.  The account concentrates on the choices that matter to a Forth
implementor, gives a complete worked example, and states the important limits
of the current native Gforth implementation.
\end{abstract}

\noindent\textbf{Keywords:} Forth; stack effect; static checking; symbolic
execution; subtyping.

\section{The problem in Forth terms}

Forth programmers already use stack comments as interfaces:
\begin{lstlisting}
: SQUARED  ( n -- n^2 )  DUP * ;
\end{lstlisting}
A useful checker must do more than count the cells on either side of
\texttt{--}.  It must know that both inputs of \texttt{*} are numeric, that
\texttt{DUP} produces two references to the same abstract value, and that a
word requiring a character address may accept an aligned address.  It must
also follow the outer interpreter: \texttt{:} consumes a name and starts a
definition, \texttt{;} ends one, comments consume source text, and immediate
control words combine non-linear source regions.

The evaluator in this project performs that work without running the program.
It takes three independent inputs:
\begin{enumerate}
  \item a finite set of types, aliases, and subtype declarations;
  \item a dictionary of word effects plus parser and control metadata;
  \item the source to be checked.
\end{enumerate}
The same engine can therefore check a small teaching vocabulary, the bundled
Forth-2012-like profile, or a system-specific profile.  The executable is
\texttt{gforth-evaluator.fs}; a Java version consumes the same files.

This is deliberately a checker for declared Forth interfaces, not a claim to
infer the behavior of arbitrary native code.  Trust begins at the primitive
specifications.  If \texttt{@} is declared incorrectly, no analysis of its
clients can repair the error.  Within that boundary, the checker infers colon
definitions, checks linear compositions, and gives source-oriented clash
diagnostics.

\section{Effects with cell identity}

\subsection{The visible suffix and its frame}

An effect is written in normal Forth order, with top of stack at the right:
\[
             (L\;--\;R).
\]
It describes only the suffix touched by the phrase.  An arbitrary deeper frame
is preserved.  Thus \texttt{( n -- n )} is applicable to \texttt{F n}, not
only to a one-cell data stack.  This implicit frame is also why effects with
different displayed depths can compose or describe two branches of one
control structure.

The empty effect \texttt{( -- )} is identity.  A clash is a separate failure
value; it is not an effect and absorbs subsequent composition.  This keeps a
reported error from being mistaken for a word that happens to have no stack
behavior.

\subsection{Types and indexed wildcards}

A profile declares canonical types and aliases.  A small address and numeric
hierarchy may include
\[
 \texttt{a-addr} < \texttt{c-addr} < \texttt{addr} < \texttt{x},
 \qquad
 \texttt{char} < \texttt{+n} < \texttt{n} < \texttt{x}.
\]
Here ``$<$'' means ``more specific than.''  The loader resolves aliases and
computes transitive relationships.  When two boundary cells are compared, the
checker keeps the more specific type if the types are comparable; unrelated
types clash.  It does not search the hierarchy for a third, common subtype.
That simple rule makes both behavior and diagnostics predictable for profile
authors.

A type name may carry an index:
\begin{lstlisting}
DUP   ( x[1] -- x[1] x[1] )
SWAP  ( x[2] x[1] -- x[1] x[2] )
OVER  ( x[2] x[1] -- x[2] x[1] x[2] )
DROP  ( x -- )
\end{lstlisting}
Equal positive indices mean equal symbolic cells.  They express the dataflow
that an ordinary unindexed effect omits.  An unindexed \texttt{x} is fresh on
each use; two plain occurrences do not assert equality.  Indices are logical
variables, not addresses and not runtime tags.  Every fetched dictionary
effect is cloned and its indices are freshened before use, so two calls of
\texttt{DUP} cannot accidentally share a variable.

\section{How a Forth phrase is composed}

Suppose the current phrase has effect $(L_1--R_1)$ and the next word has
$(L_2--R_2)$.  Composition compares $R_1$ and $L_2$ from the top of stack
inwards.  For each pair it performs the following steps:
\begin{enumerate}
  \item check that their types are comparable and select the more specific;
  \item create one fresh symbol of that type;
  \item replace both old symbols everywhere in both working effects;
  \item remove the matched output/input pair and continue.
\end{enumerate}
Replacing \emph{every} occurrence is important.  A numeric constraint found at
one boundary must reach all copies made by \texttt{DUP}, and a constraint on
one side of \texttt{SWAP} must follow the cell to its new position.  When one
touching side is exhausted, unmatched requirements or results are retained,
with the untouched stack frame inserted in the appropriate place.

For example, use these declarations:
\begin{lstlisting}
literal integer  ( -- n )
+                ( n n -- n )
DUP              ( x[1] -- x[1] x[1] )
\end{lstlisting}
The phrase \texttt{1 2 + DUP} reduces as follows:
\[
 (\,--n)\seq(\,--n)=(\,--n\;n),
\]
\[
 (\,--n\;n)\seq(n\;n--n)=(\,--n),
\]
\[
 (\,--n)\seq(x[1]--x[1]\;x[1])
       =(\,--n[1]\;n[1]).
\]
The final match specializes \texttt{x} to \texttt{n}; global substitution
applies that choice to both outputs.  The native trace makes the same process
visible:
\begin{lstlisting}
1    ( -- n )
2    ( -- n )
+    ( n n -- n[1] )
DUP  ( n[1] -- n[1] n[1] )
< n[1] n[1]
\end{lstlisting}

This procedure is deterministic.  Fresh-index numbers may differ between
runs or implementations, but consistent renaming does not alter the effect.
The older algebraic presentation relates untyped stack effects to a
polycyclic monoid \cite{poial1991,poial1994}; for using the checker, boundary
matching and whole-effect substitution are the essential facts.

\subsection{Compaction and a precision caveat}

Working effects are normalized after a successful operation.  Wildcard
numbers are made small and deterministic.  An index explicitly written by the
profile author remains visible.  An implicit identity that has too few
occurrences may be printed without its index; for example an internally
created \texttt{( n[7] -- n[7] )} can become \texttt{( n -- n )}.

Compaction produces readable inferred comments, but the implementation also
stores the compacted value.  It can therefore forget a pairwise equality that
would help a later word.  In analysis terminology this is a widening: stack
height and the type at each position remain valid, while correlation precision
can decrease.  A future version should preserve raw identities internally and
compact only when printing.

\section{Branches and control words}

Two arms of an \texttt{IF} must be represented by one contract for the code
that follows.  The checker calls the combining operation \texttt{glb}.  For
equal-depth effects it unifies corresponding inputs and outputs, again taking
the more specific of comparable types.  A mismatch in stack depth is accepted
only when the extra input cells and the same number of extra output cells form
a preserved prefix---the implicit frame made explicit by the longer arm.
Otherwise the merge clashes.

Consider two store-like branch effects requiring respectively
\[
 (\texttt{x a-addr}--), \qquad (\texttt{x c-addr}--).
\]
Because \texttt{a-addr} is more specific than \texttt{c-addr}, their common
contract requires \texttt{a-addr}.  For a complete conditional, the declared
\texttt{IF ( flag -- )} effect is composed with that branch contract, so the
selector is consumed as well.

The operation should not be read as a union of all concrete executions.  It is
a \emph{must} contract: one input/output description required to be valid for
both alternatives.  The current rule is syntactic and intentionally strict.
It catches the common Forth error where only one path leaves a value, while
allowing paths that merely expose different amounts of an otherwise preserved
frame.

Control syntax is not hard-wired into the effect algebra.  A profile can
describe a source form and how its captured regions are reduced:
\begin{lstlisting}
syntax:
  IF <then branch> [ELSE <else branch>] FI
effect:
  IF
  either <then branch> <else branch>
\end{lstlisting}
The small template language sequences control-word effects and captured
regions, combines regions with \texttt{either}, and applies \texttt{repeat} to
a loop region.  Consequently systems using \texttt{THEN}, \texttt{FI}, or a
custom immediate word can share the evaluator by changing profile data.

\section{What the loop check does---and does not do}

For a repeated region with effect $p$, the implementation computes one copy
$p$, composes two copies $p\seq p$, and asks for their common branch-style
contract:
\[
                         p^*_2 = p\meet(p\seq p).
\]
This is a finite stack-stability test.  It is useful for conventional loops
whose iteration restores the same visible stack shape, and it rejects many
bodies that grow, shrink, or change a cell incompatibly on the second
iteration.

It is not Kleene star, transitive closure, or a general abstract-interpretation
fixpoint.  Agreement of one and two iterations alone does not prove agreement
for every iteration count.  Nor does the checker discover a programmer's loop
invariant.  Profiles for \texttt{BEGIN ... UNTIL} or
\texttt{BEGIN ... WHILE ... REPEAT} must include the test word in the repeated
region when its per-iteration \texttt{flag} consumption belongs there.
Calling this rule ``idempotent repetition'' is convenient historically, but
``one-versus-two approximation'' describes the implemented guarantee more
accurately.

For production assurance, the replacement is standard: construct a
control-flow graph, maintain an abstract stack at each loop header, and
iterate to a post-fixpoint, widening if the chosen domain requires it.  The
local rule remains attractive for a small Forth tool because it terminates
immediately and works without constructing a whole-program graph.

\section{Following the outer interpreter}

A Forth source checker cannot tokenize every blank-delimited string as an
ordinary word.  The specification file therefore includes source behavior as
well as runtime stack effects:
\begin{lstlisting}
"("       parse until ")"          ( -- )
"\\"      parse until eol          ( -- )
CONSTANT  parse word define        ( x -- )
VARIABLE  parse word define        ( -- a-addr )
:         parse word define colon  ( -- )
;         control end              ( -- )
literal integer                    ( -- n )
IF        control if               ( flag -- )
\end{lstlisting}
The evaluator has interpretation and compilation states.  Defining words
consume the following name and extend the analysis dictionary; colon bodies
are checked and their inferred effects are installed at \texttt{;}.  Parser
words consume names or delimited text.  Literals are classified by profile
entries.  Word lookup is case-insensitive.

An ordinary comment such as \texttt{( n -- n )} is consumed, not trusted as an
assertion.  Otherwise a false source comment could override the very inference
being checked.  Recognized locals syntax may supply a provisional definition
shape, but that is a distinct parser feature.  Return-stack activity,
\texttt{EXECUTE}, native code generation, \texttt{THROW}, and unrestricted
metaprogramming require special abstractions or remain outside this model.

Once one colon body has been parsed, evaluation is simple.  Leaves are known
word or literal effects; adjacent nodes use composition; control nodes apply
the profile template.  Dependencies run from a node to its contained source
regions, so a non-recursive parsed body has one deterministic bottom-up result.
This modest observation is more useful than a heavy theorem here: it identifies
exactly what a reproducible diagnostic depends on---the parsed structure, the
current dictionary, and the three effect operators.

\section{Complete profile example}

The bundled \texttt{real} profile is run with
\begin{lstlisting}
./run-evaluator-gforth.sh real
\end{lstlisting}
It reads \texttt{ex1types.txt}, \texttt{ex1specs.txt}, and
\texttt{ex1prog.txt}.  The source defines:
\begin{lstlisting}
: PEEK2   DUP @ SWAP @ ;
: MAYSWAP IF SWAP FI ;
: NLOOP   BEGIN DUP 0= WHILE REPEAT ;
: ULOOP   BEGIN DUP 0= UNTIL ;
: KEEPIF  IF DUP ELSE DUP FI ;

DP DP C@ 0= KEEPIF
\end{lstlisting}
Representative profile declarations are \texttt{@ ( a-addr -- x )},
\texttt{C@ ( c-addr -- char )}, \texttt{0= ( n -- flag )}, and the indexed
stack operators shown earlier.  The inferred definitions are:
\begin{lstlisting}
PEEK2   ( a-addr[1] -- x[2] x )
MAYSWAP ( x[1] x[1] flag -- x[1] x[1] )
NLOOP   ( n[1] -- n[1] )
ULOOP   ( n[1] -- n[1] )
KEEPIF  ( x[1] flag -- x[1] x[1] )
\end{lstlisting}
The top-level phrase ends in
\begin{lstlisting}
DP      ( -- a-addr[1] )
DP      ( -- a-addr )
C@      ( a-addr -- char )
0=      ( char -- flag )
KEEPIF  ( a-addr[1] flag -- a-addr[1] a-addr[1] )
< a-addr[1] a-addr[1]
\end{lstlisting}
\texttt{C@} asks for \texttt{c-addr}; \texttt{DP} provides the more specific
\texttt{a-addr}, so the boundary succeeds.  Since \texttt{char} is below
\texttt{n}, \texttt{0=} also succeeds.  Both \texttt{KEEPIF} paths duplicate
the surviving address and their common contract retains that identity.  This
one example exercises subtyping, symbolic copies, a definition, branch merge,
loop approximation, literals, and top-level checking.

\section{Implementation and evidence}

The native implementation is one Gforth source file.  Its word families
separate type relations (\texttt{ev-ts-*}), symbols (\texttt{ev-sym-*}),
effects (\texttt{ev-spec-*}), effect lists, specification dictionaries,
source parsing, and diagnostics.  Effects are mutable records, so dictionary
entries are deep-cloned before use and independently evaluated branches receive
disjoint wildcard ranges.  Input line endings and normalized wildcard
allocation are made deterministic across runs.

Boundary matching itself is linear in the shorter touching boundary.  The
current substitution strategy scans whole working effects after each match,
so large intermediate effects can make reduction roughly quadratic.  A
union--find representation of identity classes and persistent stack rows would
reduce copying, but the straightforward implementation has proved easier to
inspect in Gforth.

The checked-in regression corpus contains 29 programs: 13 expected accepts and
16 expected rejects.  A current native batch accepts all 13 positive cases and
rejects 14 negative cases.  The two known misses are
\texttt{linear\_clash.fs} and \texttt{top\_level\_clash.fs}, both involving
the top-level phrase \texttt{DP 1+}.  This is evidence about the artifact, not
a soundness result; in particular it documents a real gap in top-level linear
checking.  The profile is also an assumption: tests cannot reveal an error
shared by a primitive declaration and its expected output.

\section{Scope and next steps}

The checker preserves the main benefits sought by a Forth programmer: inferred
stack comments, early detection of type or depth clashes, dataflow through
stack shuffles, address-type refinement, and profile-defined immediate syntax.
Its limits should remain visible:
\begin{itemize}
  \item primitive and parser specifications are trusted;
  \item only comparable types unify;
  \item compaction can lose implicit identity information;
  \item branch merge is a strict syntactic common contract;
  \item loop checking compares only one and two iterations;
  \item recursion, forward references, quotations, non-local exits, and
        unrestricted compile-time behavior are not covered generally;
  \item two known negative regression cases are currently missed.
\end{itemize}

The most direct improvements are to retain raw correlations until display, fix
the top-level clash path, and replace repeated global substitution by identity
classes.  Explicit stack-row variables would support more general
stack-polymorphic and higher-order words.  A control-flow fixpoint would give
loops a conventional invariant semantics.  Finally, checking primitive
specifications against an operational model is necessary before claiming
whole-system soundness.

Forth stack comments already provide the right interface notation.  The main
step is to make their variables, subtype relationships, parser effects, and
control combinations executable.  Symbolic stack effects do that with a small
set of understandable operations: subtype-aware boundary composition for
sequences, common contracts for alternatives, and a clearly delimited local
approximation for repetition.  The result is not a static semantics for all of
Forth, but it is a practical and extensible checker for the substantial subset
that a system profile can describe.

\begin{thebibliography}{9}
\small

\bibitem{forth2012}
Forth 200x Standardisation Committee.
\newblock \emph{Forth 2012 Standard}, 2014.
\newblock \url{https://forth-standard.org/standard/intro}.

\bibitem{poial1991}
J. Pöial.
\newblock Algebraic specifications of stack effects for Forth programs.
\newblock In \emph{1990 FORML Conference Proceedings}, pp. 282--290, 1991.

\bibitem{poial1994}
J. Pöial.
\newblock Algebraic specifications of stack effects.
\newblock \emph{Forth Dimensions}, 16(4):18--20, 1994.

\bibitem{poial2002}
J. Pöial.
\newblock Stack effect calculus with typed wildcards, polymorphism and
inheritance.
\newblock \emph{18th EuroForth Conference}, 2002.

\bibitem{poial2006}
J. Pöial.
\newblock Typing tools for typeless stack languages.
\newblock In \emph{22nd EuroForth Conference}, pp. 40--46, 2006.

\bibitem{poial2008}
J. Pöial.
\newblock Java framework for static analysis of Forth programs.
\newblock In \emph{24th EuroForth Conference}, pp. 20--23, 2008.

\bibitem{stoddart1993}
B. Stoddart and P. J. Knaggs.
\newblock Type inference in stack based languages.
\newblock \emph{Formal Aspects of Computing}, 5(4):289--298, 1993.

\end{thebibliography}

\end{document}
